\BeleskeID{08-s043}
% element: 08-s043:e04
% element: 08-s043:notes
Pre trenutka $t_0$ ulazni napon $U_0$ smatra se „beskonačno” dugo prisutnim. Intuitivno, ako bi napon na kondenzatoru bio manji od $U_0$, postojala bi struja kroz kondenzator, on bi se punio i u nekom trenutku možda bi njegov napon postao veći od $U_0$. Ako bi napon bio veći od $U_0$, struja bi tekla u suprotnom smeru, kondenzator bi se praznio i u nekom trenutku možda bi njegov napon postao manji od $U_0$. Bez obzira na smer promene, napon kondenzatora može postati jednak $U_0$. Tada struja kroz kondenzator postaje jednaka nuli i više nema promene napona.

Beskonačno vreme je idealizacija dovoljnog trajanja prethodnog stanja. Stacionarno stanje se koristi da se odredi početni napon pri kasnijem skoku, a ne da se unapred pretpostavi izlazni napon u svakom narednom intervalu.
